\section{Einleitung}\label{sec:einleitung}

Das Erfüllbarkeitsproblem der Aussagenlogik (SAT, von englisch satisfabiality "Erfüllbarkeit") ist ein 
fundamentales Problem der Informatik mit kritischen Anwendungen in Software und Hardware Überprüfung.

\subsection{Conjunctive Normal Form (CNF)}

Bei einer Boolschen Formel, die aus durch UND- und ODER-Operatoren verbunden Variablen 
besteht, (AND $\land$, OR $\lor$, NOT $\neg$) geht es beim SAT-Problem um die Frage, ob es eine Zuorndung von 
Warheitswerten zu den Variablen gibt, bei der die gesamte Formel als Wahr ausgewertet wird.

Moderne SAT-Löser arbeiten in der Regel mit Formeln in konjunktiver Normalform (CNF). Eine CNF-Formel 
besteht aus:

\begin{itemize}
    \setlength{\itemsep}{0pt}
    \setlength{\parskip}{0pt}
    \setlength{\parsep}{0pt}
    \item \textbf{Variablen}: Boolesche Variablen (z. B. $x_1$, $x_2$, ..., $x_n$)
    \item \textbf{Literale}: Eine Variable oder deren Negation (z. B. $x_1$ oder $\neg x_1$)
    \item \textbf{Klauseln}: Eine Disjunktion (ODER) von Literalen (z. B. $(x_1 \lor \neg x_2 \lor x_3)$)
    \item \textbf{Formel}: Eine Konjunktion (UND) von Satzteilen (z. B. $(x_1 \lor \neg x_2) \land (\neg x_1 \lor x_3) \land (x_2 \lor \neg x_3)$
\end{itemize}

Das CNF-Format ist sowohl ausdrucksstark (jede Boolesche Formel lässt sich in CNF umwandeln) als auch für 
algorithmisches Denken gut geeignet. Das folgende Beispiel zeigt eine einfache CNF-Datei:

\begin{verbatim}
p cnf 3 4
1 -2 0
-1 3 0
2 -3 0
-1 -2 -3 0
\end{verbatim}

Dies stellt eine Formel mit 3 Variablen und 4 Klauseln dar. Die Zeile „p cnf“ gibt den Problemtyp und die 
Dimensionen an (3 Variablen, 4 Klauseln), und jede nachfolgende Zeile stellt eine Klausel dar, die mit 0 
endet.

\newpage

\subsection{Warum SAT wichtig ist}

Trotz seiner theoretischen Bedeutung als erstes NP-vollständiges Problem (Cook, 1971) hat SAT breite 
praktische Anwendung in den Bereichen Hardware-Verifikation, Software-Testing, Planung und künstliche
Intelligenz gefunden.\cite{cook1971complexity}

Moderne SAT-solver haben in den letzten zwei Jahrzenten bemerkenswerte fortschritte gemacht; Löser wie
Kissat (Biere et al., 2020) sind in der Lage, industrielle Instanzen mit Millionen von Variablen zu
lösen.\cite{biere2020cadical}

SAT hat breite praktische Anwendung gefunden, da sich viele Probleme aus der realen Welt als 
SAT-Instanzen formulieren lassen:

\begin{itemize}
    \setlength{\itemsep}{0pt}
    \setlength{\parskip}{0pt}
    \setlength{\parsep}{0pt}
    \item \textbf{Hardware-Verifikation}: Nachweis, dass sich digitale Schaltungen korrekt verhalten
    \item \textbf{Softwaretests}: Erstellung von Testfällen, die bestimmte Pfade abdecken
    \item \textbf{Planung und Terminierung}: Ermittlung von Handlungsabläufen, mit denen Ziele erreicht werden
    \item \textbf{Kryptoanalyse}: Knacken kryptografischer Systeme durch das Lösen logischer Beschränkungen
    \item \textbf{Bioinformatik}: Analyse genetischer Sequenzen und Proteinfaltung
\end{itemize}

Die praktische Bedeutung von SAT hat jahrzehntelange Forschungen zur Entwicklung effizienter SAT-Löser 
vorangetrieben.

Der praktische Erfolg der SAT-Lösung ist weitgehend auf die bemerkenswerte Leistungsfähigkeit von 
CDCL-Lösern (Conflict-Driven Clause Learning) zurückzuführen. In der SAT-Forschung ist jedoch 
bekannt, dass kein einzelner Löser bei allen Instanzen allen anderen überlegen ist. Verschiedene 
Solver weisen komplementäre Leistungsmerkmale auf – ein Solver, der bei einer Instanzfamilie extrem 
schnell ist, kann bei einer anderen schlecht abschneiden. Diese Beobachtung hat zur Entwicklung von 
Algorithmusauswahlsystemen geführt, die automatisch den für eine bestimmte Instanz am besten 
geeigneten Solver auswählen.

Diese Erkenntnis hat zur Entwicklung des portfoliabasierten SAT-Lösens geführt, bei dem eine Sammlung 
von Solvern gepflegt wird und ein Auswahlmechanismus entscheidet, welcher Solver für eine bestimmte 
Instanz eingesetzt wird. Die zugrunde liegende Idee ist einfach, aber wirkungsvoll: Wenn wir vorhersagen 
können, welcher Solver bei einer Instanz die beste Leistung erbringt, können wir erhebliche Rechenzeit 
einsparen, indem wir Solver mit schlechter Leistung vermeiden.




